\clearpage
\begin{proofbox}
\textbf{\textcolor{CProof}{证明一：最容易理解——连续“照两次镜子”}}
\begin{center}\begin{tikzpicture}[font=\small,>=stealth,line join=round]
\draw[gray!40,rounded corners=3pt] (0,0) rectangle (14.2,8.5);\node at (7.1,7.9) {同样连续对称：两次后是否回到原高度？};
\node at (2.4,6.9) {轴轴：$f=\cos\frac{\pi x}2$};\begin{scope}[shift={(1.2666666666666666,4.348571428571429)},x=0.5666666666666667cm,y=1.0285714285714287cm]\draw[->,gray] (-1,0)--(5,0) node[right]{$x$};\draw[->,gray] (0,-1.7)--(0,1.8) node[above]{$y$};\draw[blue,thick,solid,domain=-0.7:4.8,samples=160] plot(\x,{cos(deg(pi*\x/2))});\foreach \a in {0,2}{\draw[violet,dashed] (\a,-1.4)--(\a,1.4);}\end{scope}\node[align=center,text width=4cm] at (2.4,1.7) {轴 $x=0,2$\\$d=2,\ T=4$};\node at (7.050000000000001,6.9) {点点：$g=\sin\frac{\pi x}2$};\begin{scope}[shift={(5.916666666666667,4.348571428571429)},x=0.5666666666666667cm,y=1.0285714285714287cm]\draw[->,gray] (-1,0)--(5,0) node[right]{$x$};\draw[->,gray] (0,-1.7)--(0,1.8) node[above]{$y$};\draw[teal,thick,solid,domain=-0.7:4.8,samples=160] plot(\x,{sin(deg(pi*\x/2))});\fill[orange!85!black] (0,0) circle(2pt);\fill[orange!85!black] (2,0) circle(2pt);\end{scope}\node[align=center,text width=4cm] at (7.050000000000001,1.7) {中心 $(0,0),(2,0)$\\$d=2,\ T=4$};\node at (11.7,6.9) {轴点：同一条余弦线};\begin{scope}[shift={(10.566666666666666,4.348571428571429)},x=0.5666666666666667cm,y=1.0285714285714287cm]\draw[->,gray] (-1,0)--(5,0) node[right]{$x$};\draw[->,gray] (0,-1.7)--(0,1.8) node[above]{$y$};\draw[blue,thick,solid,domain=-0.7:4.8,samples=160] plot(\x,{cos(deg(pi*\x/2))});\fill[orange!85!black] (1,0) circle(2pt);\draw[violet,dashed] (0,-1.4)--(0,1.4);\end{scope}\node[align=center,text width=4cm] at (11.7,1.7) {轴0，中心 $(1,0)$\\$d=1,\ T=4$};
\node[text width=12.6cm,align=center] at (7.1,.65) {三图同尺度：轴轴、同高点点得到2d；轴点半程翻转高度，需4d。};
\end{tikzpicture}\end{center}



先不计算函数表达式，只追踪同一个点 $(x,y)$，其中 $y=f(x)$。
\begin{itemize}[leftmargin=*,itemsep=0.4em]
\item \textbf{关于竖直直线 $x=a$ 照镜子：}$(x,y)\mapsto(2a-x,y)$。\textbf{只换左右，不换高低。}
\item \textbf{绕点 $(a,c)$ 转半圈：}$(x,y)\mapsto(2a-x,2c-y)$。\textbf{既换左右，也换高低。}
\end{itemize}

先绕横坐标为 $a$ 的位置对称一次，再绕横坐标为 $b$ 的位置对称一次。横坐标统一变为
\[
 x\to2a-x\to2b-(2a-x)=x+2(b-a).
\]
记 $h=2(b-a)$，由 $a\neq b$ 知 $h\neq0$。

\begin{description}[style=nextline,leftmargin=2.2em,itemsep=0.4em]
\item[\textbf{轴轴：}] $y\to y\to y$，于是 $f(x+h)=f(x)$，正周期取 $|h|=2d$。
\item[\textbf{点点：}] $y\to 2c-y\to2c-(2c-y)=y$，于是 $f(x+h)=f(x)$，正周期取 $|h|=2d$。
\item[\textbf{轴点：}] $y\to y\to2c-y$，第一次走 $h$ 后还没回到原高度。再走一次 $h$：$2c-y\to y$，故 $f(x+2h)=f(x)$，正周期取 $2|h|=4d$。
\end{description}

这就是\textbf{同样两次左右翻折，“高低是否翻回来”决定了周期是两倍还是四倍}。

\medskip

\end{proofbox}
\clearpage
\begin{proofbox}
\textbf{\textcolor{CProof}{证明二：考试最快——用等式一步步换进去}}
\begin{center}\begin{tikzpicture}[font=\small,>=stealth,line join=round]
\draw[gray!40,rounded corners=3pt] (0,0) rectangle (14.2,8.3);\node at (7.1,7.700000000000001) {镜子在1与3：半格出发，两次反射后到四点五};
\begin{scope}[shift={(1.6901639344262296,3.814705882352941)},x=0.9836065573770493cm,y=1.352941176470588cm]\draw[->,gray] (-0.6,0)--(5.5,0) node[right]{$x$};\draw[->,gray] (0,-1.6)--(0,1.8) node[above]{$y$};\draw[blue,thick,solid,domain=-0.2:5,samples=180] plot(\x,{cos(deg(pi*(\x-1)/2))});\fill[orange!85!black] (0.5,0.7071067811865476) circle(2.2pt);\fill[orange!85!black] (1.5,0.7071067811865476) circle(2.2pt);\fill[orange!85!black] (4.5,0.7071067811865476) circle(2.2pt);\foreach \a in {1,3}{\draw[violet,dashed] (\a,-1.3)--(\a,1.3);\node[above] at (\a,1.3){$x=\a$};}\draw[orange!85!black,dashed] (.5,.70710678)--(4.5,.70710678);\end{scope}\node[anchor=west,align=left,text width=5.0cm] at (8.4,4.0) {$f=\cos\frac{\pi(x-1)}2$\\[9pt]轴 $x=1,3$\\[9pt]$0.5\to1.5\to4.5$\\[9pt]三点同高 $\frac{\sqrt2}2$\\[9pt]横移4，两倍轴间距};
\node[text width=12.6cm,align=center] at (7.1,.65) {图中反射点来自曲线；任意起点同样满足两次反射后水平移动4。};
\end{tikzpicture}\end{center}



\textbf{先准备两条最常用的等式：}
\[
\begin{aligned}
  f(a+t)=f(a-t)&\iff f(2a-x)=f(x),\\
  f(b+t)+f(b-t)=2c&\iff f(2b-x)=2c-f(x).
\end{aligned}
\]
这里的 $t,x$ 都代表任意实数；第二条里 $b$ 可换成任意对称中心的横坐标。

\textbf{（1）轴轴：}\quad$f(2a-x)=f(x)$ 且 $f(2b-x)=f(x)$。
把第二条中的 $x$ 改成 $2a-x$：
\[
\begin{aligned}
 f\bigl(2b-(2a-x)\bigr)&=f(2a-x)\\
 f(x+2(b-a))&=f(x).
\end{aligned}
\]
故 $2d$ 是一个周期。

\textbf{（2）点点：}\quad$f(2a-x)=2c-f(x)$ 且 $f(2b-x)=2c-f(x)$。
同样把第二条中的 $x$ 改成 $2a-x$：
\[
\begin{aligned}
 f\bigl(2b-(2a-x)\bigr)
   &=2c-f(2a-x)\\
   &=2c-[2c-f(x)]=f(x).
\end{aligned}
\]
故 $2d$ 也是一个周期。\textbf{这里两处必须是同一个 $c$，否则上下翻两次不一定抵消。}

\textbf{（3）轴点：}\quad$f(2a-x)=f(x)$ 且 $f(2b-x)=2c-f(x)$。
\[
\begin{aligned}
 f\bigl(x+2(b-a)\bigr)
  &=f\bigl(2b-(2a-x)\bigr)\\
  &=2c-f(2a-x)=2c-f(x).
\end{aligned}
\]
再把 $x$ 换成 $x+2(b-a)$：
\[
\begin{aligned}
 f\bigl(x+4(b-a)\bigr)
 &=2c-f\bigl(x+2(b-a)\bigr)\\
 &=2c-[2c-f(x)]=f(x).
\end{aligned}
\]
故 $4d$ 是一个周期。\textbf{这一步必须做两次，不能看到横向移动 $2d$ 就误认为周期是 $2d$。}

若 $b<a$，式子里的平移量可能是负数，但 $f(x+h)=f(x)$ 对一切 $x$ 成立时，换元可得 $f(x-h)=f(x)$，所以可以取正的 $|h|$ 为周期。

\medskip

\end{proofbox}
\clearpage
\begin{proofbox}
\textbf{\textcolor{CProof}{证明三：与奇偶函数联系——把高度 $c$ 当作新的零线}}
\begin{center}\begin{tikzpicture}[font=\small,>=stealth,line join=round]
\draw[gray!40,rounded corners=3pt] (0,0) rectangle (14.2,8.3);\node at (7.1,7.700000000000001) {证明三：整条曲线下移3，中心高度变成0};
\begin{scope}[shift={(2.207692307692308,3.102631578947368)},x=0.9230769230769231cm,y=0.4842105263157894cm]\draw[->,gray] (-1.2,0)--(5.3,0) node[right]{$x$};\draw[->,gray] (0,-3)--(0,6.5) node[above]{$y$};\draw[blue,thick,solid,domain=-0.8:4.8,samples=180] plot(\x,{3+2*cos(deg(pi*\x/2))});\draw[orange!85!black,thick,dashed,domain=-0.8:4.8,samples=180] plot(\x,{2*cos(deg(pi*\x/2))});\fill[orange!85!black] (1,3) circle(2.2pt);\fill[orange!85!black] (1,0) circle(2.2pt);\draw[gray,dashed] (-1,3)--(5,3);\draw[<->,violet] (1.2,0)--(1.2,3);\end{scope}\node[anchor=west,align=left,text width=5.0cm] at (8.4,4.0) {蓝：$f=3+2\cos\frac{\pi x}2$\\[9pt]橙虚线：$g=f-3$\\[9pt]中心 $(1,3)\to(1,0)$\\[9pt]周期都为4\\[9pt]纵向平移不改重复距离};
\node[text width=12.6cm,align=center] at (7.1,.65) {把中线当作新的零线，关于中心的高度翻转就变成正负号交换。};
\end{tikzpicture}\end{center}



对涉及中心 $(b,c)$ 的情形，令 $g(x)=f(x)-c$。这是把图像整体向下平移 $c$，没有改变左右形状。
\[
\begin{aligned}
 f(2b-x)=2c-f(x)
 &\iff g(2b-x)=-g(x),\\
 f(2a-x)=f(x)
 &\iff g(2a-x)=g(x).
\end{aligned}
\]

\textbf{点点：}在两个中心处，$g$ 的值都\textbf{变号}；连做两次，负号抵消，$g(x+h)=g(x)$，所以 $f(x+h)=f(x)$，周期取 $2d$。

\textbf{轴点：}在对称轴处，$g$ 的值\textbf{不变号}；在中心处，$g$ 的值\textbf{变号}。连做一次轴、一次点，$g(x+h)=-g(x)$；再做一遍，$g(x+2h)=g(x)$，周期取 $4d$。

\textbf{轴轴：}两处都不变号，连做两次就有 $f(x+h)=f(x)$，周期取 $2d$。

这种看法能帮助理解：\textbf{轴对称像偶函数，点对称像奇函数；两次变号还原，一次变号还需要再来一次。}

\medskip
\textbf{\textcolor{CProof}{证明到这里，能确定的是什么？}}

我们证明了 $2d$ 或 $4d$ 确实满足周期定义，却没有证明它们是\textbf{最小的正周期}。例如曲线可能每 $\tfrac43$ 个单位就重复，但每 $4$ 个单位当然也会重复。\textbf{“一个周期”与“最小正周期”不能混写。}
\end{proofbox}
